\textbf{Proof.}
We first give explicit full-data extensions and check all boundary, support, positivity, and smoothness conditions.  Put
\[
 f_S(x)=(a+1)x^a,\qquad f_T(x)=1,\qquad 0\le x\le1.
\]
Then \(g_S(x)=a+1\) and \(g_T(x)=1\).  The declared ranges
\(c_{S,-}\le a+1\le c_{S,+}\) and
\(c_{T,-}<1<c_{T,+}\) give the required factor bounds.  Both factors are constant and hence belong to \(\mathcal H^t(M)\).  Notice also that the normalization and H\"older-radius restrictions are compatible: for any member of the class,
\[
 1=\int_0^1x^ag_S(x)\,dx
 \le \frac{\lVert g_S\rVert_\infty}{a+1}
 \le \frac{M}{a+1},
\]
so the displayed hypothesis \(M\ge a+1\) is the necessary boundary normalization, and it makes the constant choice admissible.

Keep the fixed known propensity \(e\), which obeys
\(\varepsilon\le e(x)\le1-\varepsilon\).  Conditional on \(X=x\), let the source potential outcomes be Bernoulli, conditionally independent of \(A\), with
\[
 \mu_{0,+}(x)=\mu_{0,-}(x)=\frac12,
 \qquad
 \mu_{1,\pm}(x)=\frac12\pm\delta b_h(x).
\]
Let \(A\mid X=x\) be Bernoulli with probability \(e(x)\), and impose consistency by setting \(Y=Y^S(A)\).  Give the target potential outcomes the same conditional means.  The scaling identity for a compactly supported smooth bump gives
\[
 \sup_{0<h\le1}\lVert b_h\rVert_{\mathcal H^s}
 \le C_K<\infty.
\]
Choose a fixed \(\delta>0\), independent of \(h,n,m\), sufficiently small that the two treated response functions lie in \(\mathcal H^s(M)\) and take values in \([0,1]\).  Since \(M>4\), \(\delta\) may also be chosen so that
\[
 \operatorname{Var}(Y\mid A=d,X=x)
 \ge \frac14-\delta^2\lVert K\rVert_\infty^2
 \ge M^{-1}
\]
for both arms and both signs.  Thus consistency, conditional randomization, treatment positivity, mean transportability, response smoothness, density-factor smoothness, and residual nondegeneracy all hold.  The numerical hypotheses on \(s,t\) are member properties of the fixed class and are satisfied by assumption.  Consequently the two observed laws induced by these full-data constructions belong to \(\mathcal M_{n,m}(a,s,t)\).

For \(|u|\le\delta\lVert K\rVert_\infty<1/2\), elementary Bernoulli algebra gives
\[
 \operatorname{kl}\!\left(\operatorname{Bern}\!\left(\frac12+u\right),
 \operatorname{Bern}\!\left(\frac12-u\right)\right)
 \le
 \frac{(2u)^2}{(1/2-u)(1/2+u)}
 \le C_0u^2.
\]
The source laws differ only in the treated conditional outcome.  Tensorization and \(e(x)\le1\) therefore yield
\[
\begin{aligned}
 \operatorname{KL}(P_{+,h}^{S,n},P_{-,h}^{S,n})
 &\le C_0n\delta^2\int_0^hb_h(x)^2(a+1)x^a\,dx\\
 &=C_0(a+1)\delta^2nh^{2s+a+1}
   \int_0^1K(u)^2u^a\,du\\
 &\le Cnh^{2s+a+1}.
\end{aligned}
\]
Because \(K\ge0\) is nonzero,
\(k_1:=\int_0^1K(u)\,du>0\).  The target covariate law is common to the two alternatives, and hence
\[
 \left|\tau(\mathbb P_{+,h})-\tau(\mathbb P_{-,h})\right|
 =2\delta\int_0^hb_h(x)\,dx
 =2\delta k_1h^{s+1}.
\]
This proves the source witness, including its support boundary \([0,h]\).

For the independent target-sample witness, retain
\(f_S(x)=(a+1)x^a\), put \(\ell(x)=x-1/2\), and choose a fixed sufficiently small \(\delta_0>0\).  Use the common source responses
\[
 \mu_0(x)=\frac12,
 \qquad
 \mu_1(x)=\frac12+\delta_0\ell(x),
\]
with the same Bernoulli full-data construction.  Since \(\ell\) is smooth, reducing \(\delta_0\) verifies the order-\(s\) response bounds, the outcome support, and residual nondegeneracy for every fixed \(s>0\).  Keep the entire source law fixed and set
\[
 f_{T,\pm}(x)=1\pm u_m\ell(x),
 \qquad u_m=\delta_1m^{-1/2}.
\]
The equality \(\int_0^1\ell(x)\,dx=0\) proves normalization.  The strict slack
\(c_{T,-}<1<c_{T,+}\) permits a fixed \(\delta_1>0\) such that, for every \(m\ge1\), the two densities are positive and satisfy the target factor bounds.  Reducing \(\delta_1\) if necessary also puts them in \(\mathcal H^t(M)\).  Assign target potential outcomes with the displayed common conditional means, thereby preserving mean transportability and producing admissible full-data laws.

Since \(f_{T,-}\) is uniformly bounded away from zero,
\[
\begin{aligned}
 \operatorname{KL}(P_{T,+}^{\otimes m},P_{T,-}^{\otimes m})
 &\le m\int_0^1
 \frac{\{f_{T,+}(x)-f_{T,-}(x)\}^2}{f_{T,-}(x)}\,dx\\
 &\le C_1mu_m^2\int_0^1\ell(x)^2\,dx
 \le C_2\delta_1^2.
\end{aligned}
\]
Finally, \(\int_0^1\ell(x)^2\,dx=1/12\), and hence
\[
 \left|\tau_+-\tau_-\right|
 =2\delta_0u_m\int_0^1\ell(x)^2\,dx
 =\frac{\delta_0\delta_1}{6}m^{-1/2}.
\]
The constants are independent of \(n,m,h\), and the divergence constant can be made arbitrarily small by reducing \(\delta_1\).  This proves the target witness. \(\square\)